% Chemistry Homework Formatter — student paper preamble
% Default wrapper is English (CIE papers). Load ctex only if the teacher asks for Chinese.
% Compiler: XeLaTeX. Overleaf: upload this file as main.tex plus the images/ folder.
\documentclass[10pt,a4paper]{article}

\usepackage[margin=16mm,headheight=13pt,headsep=5pt,footskip=9mm]{geometry}
\usepackage{fontspec}
\usepackage[version=4]{mhchem}
\usepackage{chemfig}
\usepackage{graphicx}
\usepackage{xcolor}
\usepackage{booktabs}
\usepackage{array}
\usepackage{enumitem}
\usepackage{needspace}
\usepackage{setspace}
\usepackage{fancyhdr}
\usepackage{amsmath,amssymb}

\raggedbottom
\setlist[enumerate]{leftmargin=1.6em,itemsep=0.25em,topsep=0.25em}
\setlist[itemize]{leftmargin=1.4em,itemsep=0.12em}

\pagestyle{fancy}
\fancyhf{}
\fancyhead[L]{\small \homeworktitle}
\fancyhead[R]{\small \homeworkmarks}
\fancyfoot[C]{\thepage}
\renewcommand{\headrulewidth}{0.3pt}

\newcommand{\homeworktitle}{Chemistry homework}
\newcommand{\homeworkmarks}{}
\newcommand{\sethomework}[2]{%
  \renewcommand{\homeworktitle}{#1}%
  \renewcommand{\homeworkmarks}{#2}%
}

\newcommand{\answerlines}[1]{%
  \par\vspace{0.3em}%
  \foreach \n in {1,...,#1}{\noindent\rule{\linewidth}{0.2pt}\\[1.05em]}%
}

% Fallback if pgffor is missing: comment \answerlines and use \vspace instead.
\usepackage{pgffor}

\graphicspath{{images/}}
\DeclareGraphicsExtensions{.png,.jpg,.jpeg,.pdf}

% Measure the image, then break only if the whole labelled clip will not fit.
\newcommand{\clipq}[3]{%
  \par
  \sbox0{\includegraphics[width=0.94\linewidth,keepaspectratio]{#3}}%
  \Needspace{\dimexpr\ht0+1.35em+2pt\relax}%
  \noindent{\small\textbf{#1}}\hfill{\small[#2]}\par\vspace{0.08em}%
  \noindent\usebox0\par
  \vspace{0.28em}%
}

\newcommand{\qfigure}[1]{%
  \sbox0{\includegraphics[width=0.94\linewidth,keepaspectratio]{#1}}%
  \Needspace{\dimexpr\ht0+0.6em\relax}%
  \noindent\usebox0\par
}

\newcommand{\qfigurebig}[1]{%
  \sbox0{\includegraphics[width=\linewidth,keepaspectratio]{#1}}%
  \Needspace{\dimexpr\ht0+0.6em\relax}%
  \noindent\usebox0\par
}

\setlength{\parindent}{0pt}
\setlength{\parskip}{0.12em}

\setlist[enumerate]{leftmargin=1.55em,itemsep=0.06em,topsep=0.12em}
\newcommand{\qhead}[1]{%
  \par\vspace{0.38em}%
  {\normalsize\bfseries #1}\par\vspace{0.12em}%
}
\newcommand{\mcq}[1]{%
  \par\noindent\textbf{#1.}\enspace\ignorespaces
}
\newcommand{\mcqend}{\hfill{[1]}\par}
\newcommand{\optrow}[4]{%
\par\noindent{\small
\begin{tabular}{@{}p{0.48\linewidth}@{}p{0.48\linewidth}@{}}
\textbf{A}\; #1 & \textbf{B}\; #2 \\[0.12em]
\textbf{C}\; #3 & \textbf{D}\; #4
\end{tabular}}\par
}
\newcommand{\optlist}[4]{%
\par
\begin{enumerate}[label=\textbf{\Alph*.}, leftmargin=1.55em, itemsep=0.05em, topsep=0.08em]
\item #1
\item #2
\item #3
\item #4
\end{enumerate}
}
\newcommand{\qfig}[2][0.58]{%
\par
\begin{center}\vspace{-0.1em}%
\includegraphics[width=#1\linewidth,keepaspectratio]{#2}%
\end{center}\vspace{-0.1em}%
}

\begin{document}
\sethomework{AS stoichiometry homework 4 · solutions and gas volumes}{16 marks}

\begin{center}
{\large\bfseries AS stoichiometry homework 4 · solutions and gas volumes}\\[0.2em]
{\small CIE 9701 AS Chemistry · 16 MCQ · 16 marks}
\end{center}

\noindent Name: \underline{\hspace{5.8cm}} \quad Class: \underline{\hspace{2.6cm}} \quad Date: \underline{\hspace{2.8cm}}

\vspace{0.25em}
{\small Topics: 13 solutions and concentration · 14 gas volumes at room conditions. Later questions also use earlier skills (mole, empirical formula, hydrates, limiting reagent).

Circle \textbf{one} letter (A--D) for each question. Show working. Calculators and a Periodic Table may be used. For gas volumes, use $24\,\mathrm{dm^3\,mol^{-1}}$ at room conditions (not $pV=nRT$), unless a question gives other conditions.}
\vspace{0.15em}

\noindent\begin{minipage}{\linewidth}
\qhead{13 · Solutions and concentration}
\mcq{1}%
Sample X is added to water and made up to a total volume of 200\,cm$^3$. This gives a solution of 0.100\,mol\,dm$^{-3}$ \ce{HCl}.

What is X?
\optlist{10\,cm$^3$ of 1.00\,mol\,dm$^{-3}$ \ce{HCl}}{30\,cm$^3$ of 0.90\,mol\,dm$^{-3}$ \ce{HCl}}{50\,cm$^3$ of 0.40\,mol\,dm$^{-3}$ \ce{HCl}}{100\,cm$^3$ of 0.30\,mol\,dm$^{-3}$ \ce{HCl}}
\mcqend
\end{minipage}\vspace{0.16em}

\noindent\begin{minipage}{\linewidth}
\mcq{2}%
Which mixture will react to form exactly one mole of water?

{\centering
{\small
\begin{tabular}{c cc}
\toprule
 & volume 2.00\,mol\,dm$^{-3}$ \ce{H2SO4}\,/\,cm$^3$ & volume 1.00\,mol\,dm$^{-3}$ \ce{NaOH}\,/\,cm$^3$ \\
\midrule
\textbf{A} & 250 & 500 \\
\textbf{B} & 250 & 1000 \\
\textbf{C} & 500 & 500 \\
\textbf{D} & 500 & 1000 \\
\bottomrule
\end{tabular}}\par}
\mcqend
\end{minipage}\vspace{0.16em}

\noindent\begin{minipage}{\linewidth}
\mcq{3}%
Citric acid is found in lemon juice.
\begin{center}citric acid\\[0.15em]\ce{HO2CCH2C(OH)(CO2H)CH2CO2H}\end{center}
Which volume of 0.40\,mol\,dm$^{-3}$ sodium hydroxide solution is required to neutralise a solution containing 0.0050\,mol of citric acid?
\optrow{12.5\,cm$^3$}{25.0\,cm$^3$}{37.5\,cm$^3$}{50.0\,cm$^3$}
\mcqend
\end{minipage}\vspace{0.16em}

\noindent\begin{minipage}{\linewidth}
\mcq{4}%
Hydrogen peroxide, \ce{H2O2}, decomposes into water and oxygen when a suitable catalyst is added.

20.0\,cm$^3$ of aqueous hydrogen peroxide decomposes to produce 600\,cm$^3$ of oxygen at room conditions.

What is the concentration of the aqueous hydrogen peroxide?
\optlist{5.00\,mol\,dm$^{-3}$}{2.50\,mol\,dm$^{-3}$}{1.25\,mol\,dm$^{-3}$}{0.625\,mol\,dm$^{-3}$}
\mcqend
\end{minipage}\vspace{0.16em}

\noindent\begin{minipage}{\linewidth}
\mcq{5}%
A washing powder contains sodium hydrogencarbonate, \ce{NaHCO3}, as one of the ingredients.

In a titration, a solution containing 1.00\,g of this washing powder requires 7.15\,cm$^3$ of 0.100\,mol\,dm$^{-3}$ sulfuric acid for complete reaction. The sodium hydrogencarbonate is the only ingredient that reacts with the acid.

What is the percentage by mass of sodium hydrogencarbonate in the washing powder?
\optrow{3.0\%}{6.0\%}{12.0\%}{24.0\%}
\mcqend
\end{minipage}\vspace{0.16em}

\noindent\begin{minipage}{\linewidth}
\mcq{6}%
The juice of one lemon reacts completely with 120\,cm$^3$ of 0.50\,mol\,dm$^{-3}$ sodium carbonate solution.

The formula of citric acid is \ce{HOOCCH2C(OH)(COOH)CH2COOH}.

No sodium carbonate is left unreacted.

What is the amount of citric acid in one lemon assuming that it is the only acid in the sample?
\optrow{0.02\,mol}{0.04\,mol}{0.06\,mol}{0.09\,mol}
\mcqend
\end{minipage}\vspace{0.16em}

\noindent\begin{minipage}{\linewidth}
\qhead{14 · Gas volumes at room conditions}
\mcq{7}%
A sample of propane, \ce{C3H8}, with a mass of 9.61\,g is completely combusted in an excess of oxygen under room conditions.

Which volume of carbon dioxide gas is produced?
\optrow{4.89\,dm$^3$}{5.24\,dm$^3$}{14.7\,dm$^3$}{15.7\,dm$^3$}
\mcqend
\end{minipage}\vspace{0.16em}

\noindent\begin{minipage}{\linewidth}
\mcq{8}%
A mixture of 10\,cm$^3$ of methane and 10\,cm$^3$ of ethane was sparked with an excess of oxygen. After cooling, the residual gas was passed through aqueous potassium hydroxide.

All gas volumes were measured at the same temperature and pressure.

Which volume of gas was absorbed by the alkali?
\optrow{15\,cm$^3$}{20\,cm$^3$}{30\,cm$^3$}{40\,cm$^3$}
\mcqend
\end{minipage}\vspace{0.16em}

\noindent\begin{minipage}{\linewidth}
\mcq{9}%
In an experiment, 0.100\,mol of propan-1-ol is burnt completely in 12.0\,dm$^3$ of oxygen, measured at room conditions.

What is the final volume of gas, measured at room conditions?
\optrow{7.20\,dm$^3$}{8.40\,dm$^3$}{16.80\,dm$^3$}{18.00\,dm$^3$}
\mcqend
\end{minipage}\vspace{0.16em}

\noindent\begin{minipage}{\linewidth}
\mcq{10}%
What is the total volume of gas produced, measured at room conditions, when 0.010\,mol of anhydrous magnesium nitrate is completely decomposed by heating?
\optrow{240\,cm$^3$}{480\,cm$^3$}{600\,cm$^3$}{720\,cm$^3$}
\mcqend
\end{minipage}\vspace{0.16em}

\noindent\begin{minipage}{\linewidth}
\mcq{11}%
The carbonate of an s-block element is reacted with an excess of hydrochloric acid.

0.833\,g of the carbonate releases 200\,cm$^3$ of gas, measured under room conditions.

What is the identity of the metal carbonate?
\optrow{\ce{Na2CO3}}{\ce{K2CO3}}{\ce{MgCO3}}{\ce{CaCO3}}
\mcqend
\end{minipage}\vspace{0.16em}

\noindent\begin{minipage}{\linewidth}
\mcq{12}%
When 0.010\,mol of a hydrocarbon X reacts with 720\,cm$^3$ of hydrogen at room conditions, an alkane is formed.

What is hydrocarbon X?
\qfig[0.72]{fig-q12-structures.png}
\mcqend
\end{minipage}\vspace{0.16em}

\noindent\begin{minipage}{\linewidth}
\qhead{Earlier skills · mole, empirical, hydrate, limiting}
\mcq{13}%
How many molecules are present in 62\,g of solid white phosphorus, \ce{P4}?
\optrow{$L$}{$2L$}{$\dfrac{L}{2}$}{$\dfrac{L}{4}$}
\mcqend
\end{minipage}\vspace{0.16em}

\noindent\begin{minipage}{\linewidth}
\mcq{14}%
Hexamine is a crystalline solid used as a fuel in portable stoves.

The diagram shows its skeletal structure.
\qfig[0.28]{fig-q14-hexamine.png}
What is the empirical formula of hexamine?
\optrow{\ce{CH2N}}{\ce{C3H6N2}}{\ce{C4H8N4}}{\ce{C6H12N4}}
\mcqend
\end{minipage}\vspace{0.16em}

\noindent\begin{minipage}{\linewidth}
\mcq{15}%
Barium dithionate, \ce{BaS2O6.2H2O}, is soluble in water.

\ce{S2O6^2-} ions slowly decompose in acidic solution.
\[\ce{S2O6^{2-}(aq) -> SO2(g) + SO4^{2-}(aq)}\]
3.513\,g of \ce{BaS2O6.2H2O} is dissolved in water in a 100\,cm$^3$ volumetric flask and the solution made up to the mark with \ce{HCl(aq)}.

At time $x$\,min, a white precipitate of mass 0.661\,g is present in the flask.

What is the concentration of \ce{BaS2O6} in the volumetric flask at time $x$\,min?
\optrow{0.0077\,mol\,dm$^{-3}$}{0.0090\,mol\,dm$^{-3}$}{0.077\,mol\,dm$^{-3}$}{0.090\,mol\,dm$^{-3}$}
\mcqend
\end{minipage}\vspace{0.16em}

\noindent\begin{minipage}{\linewidth}
\mcq{16}%
An excess of chlorine was bubbled into 100\,cm$^3$ of hot 6.0\,mol\,dm$^{-3}$ sodium hydroxide.

How many moles of sodium chloride would be produced in the reaction?
\optrow{0.30}{0.50}{0.60}{0.72}
\mcqend
\end{minipage}\vspace{0.16em}

\end{document}
